\section{Background and related work}
\textbf{BioCLIP.} BioCLIP~\citep{stevens2024bioclip} is a CLIP ViT-B/16 trained contrastively on TreeOfLife-10M,
10M images from 454k taxa, with text that mixes the full Linnaean taxonomy and common names. Its public demo
classifies an image against the text embeddings of every species of TreeOfLife-10M. BioCLIP-2~\citep{gu2025bioclip2}
scales the recipe to 214M images.

\textbf{Adversarial examples beyond $\ell_p$ balls.} Generative models produce adversarial examples that are not
small perturbations: unrestricted examples from GANs~\citep{song2018constructing}, naturally occurring hard examples
(ImageNet-A, \citealp{hendrycks2021natural}), and diffusion-based attacks that edit images in latent space. We target a
scientific model and require fidelity to be verified independently.

\textbf{Aligning diffusion models with a reward.} DDPO~\citep{black2024ddpo} casts denoising as a multi-step
decision process and optimises a reward with policy gradients; GRPO~\citep{shao2024deepseekmath} normalises rewards
within a group of samples of the same prompt. RAFT~\citep{dong2023raft} fine-tunes on the best samples ranked by the
reward. Optimising a learned reward invites reward hacking~\citep{skalse2022defining}.

\textbf{Fast sampling.} DDIM~\citep{song2021ddim} and DPM-Solver++~\citep{lu2022dpmsolverpp} reduce the number of
denoising steps; latent consistency models~\citep{luo2023lcmlora} sample in 2--4 steps; latent diffusion
\citep{rombach2022ldm} moves the computation to a compressed latent space, and tiny decoders (TAESD) make decoding
cheap. SDEdit~\citep{meng2022sdedit} edits a real image by noising it and denoising it with a new prompt.
